\chapter{总结与展望}

\section{工作总结}

本文围绕企业供应链协同管理场景，完成了基于 Spring Boot 与 Vue 3 的企业供应链协同管理系统的分析、设计与实现工作。针对传统管理模式下存在的信息分散、流程衔接不畅、库存状态不透明和质量追溯链条不完整等问题，本文从业务协同需求出发，设计了覆盖顾客、销售管理员、仓库管理员、生产管理员、采购管理员、供应商、质检管理员和系统管理员等多角色的统一业务平台。

在系统设计方面，本文完成了需求分析、总体架构设计、权限控制设计、业务流程设计和数据库设计，明确了订单驱动下的销售审核、库存处理、生产执行、采购供应、质检留痕和质量追溯等关键业务链路。在系统实现方面，本文基于 Spring Boot 构建后端服务，结合 Spring Security 与 JWT 实现认证授权，采用 Vue 3、Vite 与 Axios 构建前端交互界面，以 MySQL 和 JPA 管理业务数据，并利用 WebSocket/STOMP 实现实时消息推送。

在核心功能层面，本文重点实现了基于角色的访问控制机制、订单—仓储—生产—质检协同流程、采购—供应商—仓库协同流程、库存预警与周计划生成机制、实时消息同步机制以及质量追溯功能。测试结果表明，系统能够较好地支撑企业供应链业务的基本运行需求，在流程贯通性、数据一致性和业务可追溯性方面具有较好的应用价值。

\section{研究不足与未来展望}

尽管本文完成了系统的主要设计与实现工作，但仍存在一定不足。首先，周计划生成目前主要基于规则和阈值计算，对复杂需求波动、季节性变化和多约束排程问题的适应能力仍较弱。其次，系统当前更偏重业务流程协同，在经营分析、数据可视化和决策辅助方面的深度仍有提升空间。再次，本文实现主要面向 Web 端，尚未针对移动端场景进行专门优化，在现场作业和移动审批等方面仍不够完善。此外，当前系统采用单体架构实现，虽然便于开发与维护，但在更大规模业务场景下的弹性伸缩和模块解耦能力仍有待增强。

后续研究可从以下几个方向展开。其一，可引入更精细的需求预测方法和多因素计划优化策略，增强生产与采购计划生成的准确性与适应性。其二，可进一步完善可视化分析功能，围绕订单履约、库存周转、供应商响应和质量问题分布等内容构建管理看板。其三，可扩展移动端应用，提高仓储、质检和采购审批等场景下的便捷性。其四，可探索系统的微服务化改造，增强模块独立部署能力和高并发场景下的系统稳定性。由此可见，企业供应链协同管理系统仍具有较大的优化空间和工程应用前景。